% =====================================================================
% Figure: Functorial commutative square for TCDA (Section 3.3 / 3.5)
% Working draft -- the user will refine final visual styling.
% =====================================================================
\begin{figure}[t]
\centering
\begin{tikzcd}[column sep=large, row sep=large]
P^{X = x}
  \arrow[r, "\Tc"] \arrow[d, "\doop(x') / \doop(x)"', shift right=2pt]
  \arrow[d, "T_{\!\Cc}"', shift left=10pt, dashed]
  & \Tc(P^{X=x})
  \arrow[d, "\Phi_{\Cc}", shift left=2pt]
  \\
P^{X = x'}
  \arrow[r, "\Tc"']
  & \Tc(P^{X=x'})
\end{tikzcd}
\caption{Functorial picture of a TCDA query. The horizontal arrows apply
a topological functor $\Tc$ (e.g.\ $\PHk{k}$, $\Reeb$, $\Map$) to two
interventional distributions. The left vertical arrow $\doop(x)\!\to\!\doop(x')$
acts at the level of the causal model; the right vertical arrow $\Phi_{\Cc}$
acts at the level of topological summaries. A TCDA solution requires either
that the diagram \emph{commutes} -- so that the topological summary of a
causal contrast equals the causal contrast of topological summaries -- or
that the failure to commute is itself a quantifiable, interpretable causal
object. Stability of $\Tc$ (cf.\ \Cref{eq:stability}) makes the right
vertical arrow Lipschitz with respect to bottleneck distance; hence under
commuting the left vertical arrow inherits a Lipschitz bound in the same
metric.}
\label{fig:functorial}
\end{figure}
